\documentclass[11pt,a4paper]{article}

% --- Кодировки и язык ---
\usepackage[utf8]{inputenc}
\usepackage[T2A]{fontenc}
\usepackage[russian]{babel}

% --- Геометрия и шрифт ---
\usepackage[top=2cm,bottom=2cm,left=2.2cm,right=2.2cm]{geometry}
\usepackage{microtype}
\usepackage{parskip} % абзацы без отступа, с интервалом — читается легче в коротком тексте
\usepackage{amsmath}
\usepackage{amssymb}

% --- Математика ---
\usepackage{amsmath,amssymb,amsthm}
\newtheorem{lemma}{Лемма}
\newtheorem{theorem}{Теорема}
\newtheorem{proposition}{Утверждение}
\newtheorem{definition}{Определение}
\theoremstyle{remark}
\newtheorem*{remark}{Замечание}


% --- Ссылки и прочее ---
\usepackage[colorlinks=true,linkcolor=blue,urlcolor=blue,citecolor=blue]{hyperref}
\usepackage{enumitem}
\setlist{nosep,leftmargin=1.4em}

% --- Колонтитулы ---
\usepackage{fancyhdr}
\pagestyle{fancy}
\fancyhf{}
\fancyhead[L]{\small Разбор статьи: Direct Preference Optimization}
\fancyhead[R]{\small \thepage}
\renewcommand{\headrulewidth}{0.4pt}

% --- Метаданные ---
\title{\vspace{-1.5cm}\textbf{Анализ статьи «Direct Preference Optimization: Your Language Model is Secretly a Reward Model»}}
\author{Зарубин Егор}
\date{\today}

\begin{document}
\maketitle
\vspace{-0.8cm}

% =========================================================
\section{Анализ и оценка статьи}
% =========================================================

\subsection{Задача и формализация}
Для начала зафиксируем, что предлагают авторы статьи: 
\begin{enumerate}
    \itemПредлагается в обычном пайплайне RLHF, используемом в предыдущих работах и состоящем из трех этапов, заменить обучение отдельной модели награды и последующую RL-оптимизацию политики одним этапом обучения.
    \item Утверждается, что DPO стабильный производительный и вычислительно легкий.
    \item Дообучение с помощью DPO лучше RLHF контролирует тональность генераций и улучшает качество ответов при суммировании текста и одноходовом диалоге.
    \item DPO более прост в реализации и обучении, чем RLHF.
\end{enumerate}

Таким образом, авторы рассматривают уже существующую задачу RLHF, то есть задачу максимизации ожидаемой награды со штрафом в виде KL-дивергенции относительно исходной политики. Их цель - это предложить более простой способ решения этой задачи, позволяющий обучать модель непосредственно на данных о человеческих предпочтениях.

В статье используется следующая постановка: фиксируем референсную модель $\pi_{\text{ref}}$ и модель предпочтений 
$$p^*(y_1 \succ y_2 \mid x) = \frac{\exp (r^*(x, y_1))}{\exp (r^*(x, y_1)) + \exp (r^*(x, y_2))}.$$
Для обучения итоговой политики решаем оптимизационную задачу:
$$\max_{\pi_\theta} \mathbb{E}_{x \sim \mathcal{D}, y \sim \pi_\theta(y \mid x)} [r_\phi(x, y)] - \beta \mathbb{D}_{\text{KL}} [\pi_\theta(y \mid x) \parallel \pi_{\text{ref}}(y \mid x)].$$
Авторы используют стандартную постановку RLHF, рассмотренную в предыдущих работах. При этом для дальнейшего вывода DPO существенны как выбранная KL-регуляризация, так и предположение о модели человеческих предпочтений. 

Хотя постановка и соответствует цели авторов, поскольку они предлагают новый способ решения уже существующей задачи, полученные теоретические результаты зависят от выбранной модели предпочтений и KL-регуляризации, что ограничивает их применимость к другим постановкам RLHF. Кроме того, авторы рассматривают конкретный, распространенный в предыдущих работах RLHF-пайплайн, поэтому результаты экспериментального сравнения нельзя автоматически переносить на другие его варианты.

\subsection{Теоретическая идея DPO и ее оценка}

Что по факту делают авторы: они берут уже известную формулу для оптимальной политики при известной функции награды
\[\pi_r(y \mid x) = \frac{1}{Z(x)} \pi_{\text{ref}}(y \mid x) \exp \left( \frac{1}{\beta} r(x, y) \right),\]
выражают из нее функцию награды и подставляют ее в функцию потерь для модели предпочтений и получают новый лосс, который зависит только от предпочтений, референсной и обучаемой политик, но не от функции награды. Полученную функцию называют \textit{DPO-лоссом.}

При алализе новой функции потерь они делают большой акцент на члене с сигмоидой
\[\begin{aligned}
& \nabla_\theta \mathcal{L}_{\text{DPO}}(\pi_\theta; \pi_{\text{ref}}) = \\
& \quad - \beta \mathbb{E}_{(x, y_w, y_l) \sim \mathcal{D}} \left[ \underbrace{\sigma(\hat{r}_\theta(x, y_l) - \hat{r}_\theta(x, y_w))}_{\text{больший вес, когда оценка награды неверна}} \left[ \underbrace{\nabla_\theta \log \pi(y_w \mid x)}_{\text{увеличить вероятность } y_w} - \underbrace{\nabla_\theta \log \pi(y_l \mid x)}_{\text{уменьшить вероятность } y_l} \right] \right],
\end{aligned}\]
хотя соответствующий коэффициент автоматически возникает при вычислении градиента предложенной функции потерь. Несмотря на это, авторы пытаются убрать этот член и утверждают, что это приводит к "дегенерации модели". В доказательство этому приводится 2 примера, когда модель зацикливается при генерации ответа. Однако такие примеры лишь демонстрируют возможность деградации, но не позволяют судить о частоте ее возникновения.

Также в приложении A.4 при выводе градиента допущены ошибки в знаках, однако в основной части статьи приведена верная формула.

В пятой части и в теореме 1 в частности авторы доказывают, что таким преобразованием они не сузили класс представимых функций награды с точностью до эквивалентности, однако она не гарантирует восстановления оптимальной политики при минимизации DPO-лоса. Для этого необходимы дополнительные предположения о данных, выразительности семейства политик и свойствах оптимизации.
Все вспомогательные утверждения и сама теорема доказаны верно, и, как утверждают сами авторы, доказательства тривиальны.
 
В \S 5.2 авторы обращают внимание на нормировочный член $\beta\log Z(x)$, который зависит только от запроса $x$, но не от ответа $y$, поэтому его вычитание из функции награды не меняет оптимальную политику. Однако, при стохастическом вычислении градиента такое вычитание может уменьшить дисперсию его оценки, тем самым повысив стабильность обучения. В обычном RLHF вычисление этого члена затруднительно, поскольку требует суммирования по всем возможным ответам, тогда как в DPO благодаря выбранной параметризации он равен нулю и не требует вычисления.
Однако, гипотеза о влияении члена на обучение не проверяется в разделе экспериментов и авторы не возращаются к этой идее в обсуждении, поэтому предложенное объяснение нестабильности не получает отдельного экспериментального подтверждения.

Итого, получаем, что основная идея DPO основана на относительно простом преобразовании уже известных результатов. Математический переход от модели награды к политике обоснован, но объяснение практической стабильности DPO существенно слабее, чем обоснование его функции потерь.

\subsection{Экспериментальная часть и ее оценка}
Для экспериментов были выбраны 3 задачи
\begin{enumerate}
    \item Продолжение рецензии фильма с нужной тональностью, здесь DPO сравнивается с preffered-finetuned моделью, полным пайплайном PPO и моделью, которая обучена на максимизировать вероятность желаемых ответов и минимизировать нежелаемых(Unlikelihood)
    \item Суммаризация текста, здесь также сравнение c GPT-J, ее SFT, PFT, а также в BEST-N SFT, то есть берем N раз ответ SFT и выбираем лучший.
    \item Одноходовый диалог, здесь аналогичное второй задаче сравнение с Pythia-2.8B
\end{enumerate}
Первый эксперимент позволяет проверить, насколько хорошо DPO оптимизирует компромисс между наградой и KL-дивергенцией, второй и третий, насколько хорошо DPO работает на реальных данных человеческих предпочтений.
Сравнение с best-N приведено, чтобы проверить, возможно ли добиться того же результата улучшив качество на этапе генерации, а не на этапе обучения.

Набор задач позволяет проверить то, что заявляют авторы, однако проведенных экспериментов недостаточно для этого. Как мы заметели в начале обзора, было заявлено, что DPO стабильный производительный и вычислительно легкий. Несмотря на то, что авторы демонстрируют преимущество DPO по качеству в рассмотренных задачах, заявленные стабильность и вычислительная эффективность не получают столь же подробной экспериментальной проверки. Хотя этого и не было заявлено, но не было проведено экспериментов, которые бы сравнивали модели при равном бюджете ресурсов. Также, эксперименты не показывают, насколько результаты сохраняются при изменении размера модели, ее архитектуры, датасета предпочтений и RLHF-пайплайна.

Кроме того, вопросы вызывает выбор методов для сравнения: авторы рассматривают DPO в сравнении с конкретным пайплайном RLHF на основе PPO и несколькими простыми методами обучения, однако не сравнивают его с другими существовавшими на тот момент подходами, которые также позволяют обучать модель непосредственно на данных о предпочтениях, например RRHF~\cite{rrhf} и SLiC-HF~\cite{slic-hf}. Это ограничивает возможность оценить преимущества DPO относительно альтернативных методов.

Отдельные вопросы вызывает экспериментальная часть \S 6.2. Во-первых, авторы утверждают, что DPO более устойчив к изменению температуры генерации, чем PPO. Однако, не обязательно такая устойчивость обязательно является преимуществом, ведь температура используется для управления случайностью генерации, независимость качества от температуры недостатком, так как может сопровождаться меньшим разнообразием ответов. Авторы оценивают только долю побед модели, но не исследуют, как изменение температуры влияет на разнообразие генераций. Поэтому сделать вывод о том, что такая устойчивость является безусловным преимуществом DPO, нельзя.

Также авторы признают, что не проводили существенной настройки гиперпараметра $\beta$ для DPO, предполагая, что это могло привести к недооценке качества их метода. Однако $\beta$ непосредственно влияет на функцию потерь DPO, поэтому без дополнительной настройки нельзя оценить, насколько близки полученные результаты к наилучшим возможным. С другой стороны, из отсутствия настройки не следует, что качество DPO действительно могло быть существенно улучшено. Для более объективного сравнения стоило бы исследовать зависимость качества обоих методов от гиперпараметров при сопоставимых вычислительных затратах на их подбор.

Наиболее существенные вопросы вызывает эксперимент с Anthropic HH. В отличие от задачи суммаризации, здесь авторы не проводят полноценного сравнения DPO и PPO, обученных на одной исходной модели. Вместо этого они используют уже обученную PPO-модель, для которой не смогли подобрать параметры генерации, позволяющие превзойти исходную Pythia-2.8B. После этого авторы предлагают использовать Best-of-128 как приблизительный ориентир качества PPO, опираясь на результаты предыдущего эксперимента. Однако Best-of-N и PPO являются принципиально разными подходами: первый выбирает лучший ответ из нескольких сгенерированных, тогда как второй изменяет саму политику. Поэтому результаты этого эксперимента показывают преимущество DPO над рассмотренными методами, но не позволяют судить о его превосходстве над хорошо настроенным PPO в задаче одноходового диалога.

В качестве основного оценщика качества ответов в задачах суммаризации и одноходового диалога авторы используют GPT-4. Ссылаясь на предыдущие исследования, они говорят, что языковые модели могут лучше оценивать качество текста, чем автоматические метрики. Чтобы дополнительно обосновать использование GPT-4, в \S 6.3 авторы сравнивают его оценки с человеческими.

Однако, достоверность этого эксперимента вызывает вопросы. В нем приняли участие всего 25 студентов и выпускникоы Стэнфорда. Такую выборку трудно назвать репрезентативной. Кроме того, исследование проводится только на задаче суммаризации и нескольких заранее выбранных сравнениях моделей, поэтому неясно, насколько результаты можно обобщить на остальные задачи.

В таблице 1 авторы сравнивают доли побед по оценкам GPT-4 и людей и частоту совпадения их оценок. На основании этих результатов делается вывод, что GPT-4 согласуется с людьми примерно так же часто, как люди друг с другом. Но по приведенным результатам нельзя однозначно сделать такой вывод при отсутсвии доверительных интервалов.

Таким образом, эксперимент дает слабые основания использовать GPT-4 для автоматической оценки. И поскольку значительная часть результатов статьи основана именно на его оценках, это ограничение существенно для интерпретации экспериментальных выводов.

\subsection{Новизна и нетривиальность результатов}

Основная идея DPO основана на уже известных математических результатах. В частности, формула оптимальной политики при KL-дивергенции была получена задолго до DPO и использовалась, например, в AWR \cite{awr}. Более того, сама идея объединения обучения модели награды и политики не нова: в ASAF \cite{asaf} похожий подход используется для обучения политики через классификатор. Также до DPO существовали методы непосредственного обучения на человеческих предпочтениях без отдельной модели награды и RL-оптимизации, например Direct Preference-based Policy Optimization \cite{an2023}, RRHF \cite{rrhf} и SLiC-HF \cite{slic-hf}. Таким образом, новизна DPO заключается не в отдельных математических инструментах или самой идее обучения на предпочтениях, а в их сочетании: авторы используют аналитическую связь между наградой и оптимальной политикой, чтобы получить функцию потерь непосредственно из модели Bradley–Terry. На мой взгляд, несмотря на относительную простоту математических преобразований, предложенный подход представляет собой нетривиальное применение известных результатов, позволяющее существенно упростить стандартный RLHF-пайплайн.

% =========================================================
\section{Идеи в продолжение статьи}
% =========================================================

\subsection{Обобщение модели предпочтений}

В DPO используется модель предпочтений Bradley–Terry, однако можно рассмотреть более общий случай. Предположим, что человек оценивает каждый ответ с некоторой случайной ошибкой $U(x,y)=r(x,y)+\varepsilon_y,$
где \(\varepsilon_y\) — норсв с произвольным непрерывным распределением.
Тогда вероятность предпочтения одного ответа другому имеет вид
\[
P(y_1\succ y_2\mid x)
=P(\varepsilon_2-\varepsilon_1<r_1-r_2)
=F(r_1-r_2),
\]
где \(F\) — функция распределения разности ошибок.

Используя ту же связь между наградой и оптимальной политикой, что и в DPO, получаем обобщенную функцию потерь:
\[
\mathcal L_F(\theta)
=-\mathbb E_{\mathcal D}\log F(\beta\log\frac{\pi_\theta(y_w\mid x)}
{\pi_{\mathrm{ref}}(y_w\mid x)}
-\beta\log\frac{\pi_\theta(y_l\mid x)}
{\pi_{\mathrm{ref}}(y_l\mid x)}),
\]

Таким образом, стандартный DPO получается как частный случай при \(F=\sigma\) и такое обобщение не меняет выкладок. При этом для дифференцируемой \(F\) вес примера в градиенте равен \(F'(\Delta_\theta)/F(\Delta_\theta)\). Например, для сигмоиды, используемой в DPO, вес равен $\sigma(-\Delta_\theta)$ и ограничен единицей. Если же предположить, что ошибки имеют нормальное распределение, то $F=\Phi$, вес равен $\varphi(\Delta_\theta)/\Phi(\Delta_\theta)$ и неограниченно возрастает при $\Delta_\theta\to-\infty$. Таким образом, выбор модели предпочтений существенно влияет на вклад примеров, для которых предсказание противоречит наблюдаемому предпочтению. Было бы интересно исследовать, как это сказывается на устойчивости обучения.

Также интересно рассмотреть обобщение в другом направлении: можно ли заменить Кульбака-Лейблера на другую дивергенцию? Ведь возможно другой регуляризатор позволит получать более разнообразные генерации, для этого будет полезно исследовать, какие условия нужно наложить на дивергенцию, чтобы можно было выразить функцию награды через оптимальную политику.

\subsection{Дополнительные эксперименты}

В первую очередь хотелось бы более подробно проверить заявленные авторами стабильность и вычислительную эффективность DPO. Для этого можно провести сравнение DPO, PPO и других похожих алгоритмов, например RRHF и SLiC-HF, при сопоставимых вычислительных ресурсах. Предлагается взять одинаковые исходные модели, датасеты и оборудование, выделить каждому алгоритму одинаковый бюджет GPU-часов, включая подбор гиперпараметров, и сравнить достигнутое качество. Отдельно интересно исследовать эффективность базовых реализаций алгоритмов без ограничения бюджета: сравнить время обучения, количество эпох до достижения заданного качетсва, кривую лоса. Также хотелось бы экспериментально проверить гипотезу из \S 5.2 о влиянии нормировочного члена на дисперсию градиента. Для этого можно сравнить обучение PPO с использованием соответствующей нормализации награды и без нее, измеряя дисперсию оценок градиента и скорость сходимости. Наконец, поскольку авторы не проводили существенной настройки параметра $\beta$, стоило бы исследовать зависимость качества DPO, отклонения от референсной политики и стабильности обучения от его значения. Также хотелось бы проверить, не является ли заявленная авторами устойчивость DPO к изменению температуры генерации следствием меньшей гибкости модели. Для проверки этой гипотезы можно сравнить, как меняются качество и разнообразие генераций DPO и PPO при разных температурах. Такие эксперименты позволили бы более объективно оценить преимущества DPO и выяснить, в каких условиях они сохраняются. 

\section{Выводы}

Эксперименты демонстрируют хорошие результаты DPO в рассмотренных задачах, однако не все заявления авторов получают одинаково убедительное подтверждение. В частности, остаются вопросы относительно стабильности обучения, вычислительной эффективности и преимуществ DPO перед другими методами обучения на предпочтениях.

Несмотря на эти ограничения, я считаю DPO значимым результатом, прежде всего благодаря простоте предложенного подхода. При этом как теоретические обобщения метода, так и более подробное исследование его практических свойств представляют интерес для дальнейшей работы.

\begin{thebibliography}{9}

\bibitem{rrhf}
Z. Yuan et al.
\textit{RRHF: Rank Responses to Align Language Models
with Human Feedback without tears}.
arXiv:2304.05302, 2023.
\url{https://arxiv.org/abs/2304.05302}

\bibitem{slic-hf}
Y. Zhao, R. Joshi, T. Liu, M. Khalman, M. Saleh, P. J. Liu.
\textit{SLiC-HF: Sequence Likelihood Calibration
with Human Feedback}.
arXiv preprint arXiv:2305.10425, 2023.
\url{https://arxiv.org/abs/2305.10425}

\bibitem{awr}
X. B. Peng, A. Kumar, G. Zhang, S. Levine.
\textit{Advantage-Weighted Regression:
Simple and Scalable Off-Policy Reinforcement Learning}.
arXiv:1910.00177, 2019.
\url{https://arxiv.org/abs/1910.00177}

\bibitem{asaf}
P. Barde et al.
\textit{Adversarial Soft Advantage Fitting:
Imitation Learning without Policy Optimization}.
NeurIPS, 2020.
\url{https://arxiv.org/abs/2006.13258}

\bibitem{an2023}
G. An et al.
\textit{Direct Preference-based Policy Optimization
without Reward Modeling}.
arXiv:2301.12842, 2023.
\url{https://arxiv.org/abs/2301.12842}

\end{thebibliography}

\end{document}